\documentclass[11pt,a4paper]{article}

% Packages
\usepackage[utf8]{inputenc}
\usepackage{amsmath}
\usepackage{graphicx}
\usepackage{booktabs}
\usepackage{cite}
\usepackage{hyperref}

% =============================
% Title
% =============================
\title{A Deep Learning Model for Abnormal Human Activity Recognition}
\author{Sanele Hlabisa \\
University of KwaZulu-Natal \\
\texttt{218023135@stu.ukzn.ac.za}}
\date{}

\begin{document}

\maketitle

% =============================
% Abstract
% =============================
\begin{abstract}

\end{abstract}

% =============================
% Introduction
% =============================
\section{Introduction}

\subsection{Background}

\subsection{Motivation}

\subsection{Contributions}

% =============================
% Related Work
% =============================
\section{Related Work}

\subsection{Traditional Approaches}

\subsection{Deep Learning Approaches}

\subsection{Summary of Gaps}
% What is missing in literature?

% =============================
% Proposed Method
% =============================
\section{Proposed Method}

\subsection{Problem Formulation}
% Define problem mathematically if needed.

\subsection{Model Architecture}
% Describe architecture.

\subsection{Training Procedure}
% Loss function, optimization, hyperparameters.

\subsection{Implementation Details}
% Hardware, dataset split, preprocessing.


% =============================
% Experimental Results
% =============================
\section{Experimental Results}

\subsection{Dataset Description}
% Describe dataset.
During this resarch we used the 

Abnormal Activities 11 classes ~1070 videos
Classes
    Knife Hazard:
    Robery:
    Pollution:
    Harrasment:
    Property Damage:
    Hijack:
    Terrorism:
    Begging:
    Normal Videos:
    Drunkeness:

AIRTLab has 2 classes Violence and Non Violence: and ~350 videos, recoded with camera fromdifferen tviews
Classes:
    Violence:
    Non-Vidoence:

Train test split used for this experiment was 70:15:15 for train, validation and test dataset

\subsection{Evaluation Metrics}
% Accuracy, F1, etc.
Accuracy, Precision, Recall, F1 Score and Confusion Matrixs where choosen as perforamce indicators for the proposed ConvLSTM model

\subsection{Quantitative Results}
% Tables and comparisons.

\subsection{Ablation Study}
% Optional.


% =============================
% Discussion
% =============================
\section{Discussion}

\subsection{Analysis of Results}
% Interpret results.

\subsection{Limitations}
% Weaknesses.

% =============================
% Conclusion and Future Work
% =============================
\section{Conclusion and Future Work}

\subsection{Conclusion}
% Summarize findings.

\subsection{Future Work}
% Next steps.

% References

\bibliographystyle{IEEEtran}
\bibliography{references}

\end{document}